\chapter{Einleitung}

\section*{Motivation und Kontext}

Wir leben in einer Zeit rasant wachsender Urbanisierung, zunehmender Klimavariabilität und eines steigenden Bewusstseins für Luftqualität als Faktor der öffentlichen Gesundheit. Dennoch bleiben detaillierte, hyperlokal aufgelöste Umweltdaten überraschend knapp. Die meisten Wetter- und Luftqualitätsstationen sind kostspielig, stromhungrig und an nur wenigen offiziellen Standorten pro Stadt installiert. Die Lücken zwischen diesen Stationen können Kilometer umfassen, sodass ganze Stadtviertel ohne aussagekräftige Daten über ihre lokale Luftqualität bleiben.

Das Problem der Datenlücken ist jedoch nicht auf urbane Räume beschränkt. In ländlichen Gebieten, auf Versuchsflächen oder an abgelegenen Messstationen fehlen oft grundlegende Infrastrukturen wie WLAN oder Mobilfunkversorgung. Wo Mobilfunk verfügbar ist, erweist sich dieser häufig als wirtschaftlich unrentabel oder energetisch prohibitiv, insbesondere für Langzeitmessungen in autonomen Systemen. Herkömmliche funknetzbasierte Sensoren benötigen erhebliche Energiemengen und erzeugen damit hohe Betriebskosten.

Gleichzeitig hat sich die IoT-Landschaft enorm weiterentwickelt. Low-Power-Wide-Area-Netzwerke wie LoRaWAN ermöglichen es heute, Sensordaten über mehrere Kilometer mit nur einem Bruchteil der von Wi-Fi oder Mobilfunk benötigten Energie zu übertragen, oft im Mikrowatt-Bereich im Ruhemodus. In Kombination mit kostengünstiger Solar-Energiegewinnung und modernen Ultra-Low-Power-Mikrocontrollern wird eine vollständig autarke, unbeaufsichtigte Messstation nicht nur machbar, sondern auch an Orten ohne externe Energiequellen oder Telekommunikationsinfrastruktur praktisch und wirtschaftlich.

\section*{Projektrahmen}

Das Projekt \textit{HydroNodeStation01} entstand mit dem Anspruch, eine Station zu schaffen, die kosteneffizient reproduzierbar ist und ohne regelmäßige manuelle Wartung betrieben werden kann. Kernziel ist die Langzeitüberwachung an entlegenen oder schwer zugänglichen Standorten, wo traditionelle Infrastruktur fehlt oder unwirtschaftlich ist. Die Station soll Wochen bis Monate vollständig autark durch Solarenergie versorgt werden, ohne dass der Betreiber eingreifen muss. Dadurch werden Messungen auch an entlegenen Standorten praktikabel, an denen regelmäßige Wartungsfahrten kostspielig oder logistisch schwierig wären.

Gleichzeitig ist das Design so offen gestaltet, dass jeder die Hardware reproduzieren, modifizieren und in ein Datennetzwerk integrieren kann. Auf diese Weise entsteht eine Grundlage für dezentralisierte, bürgerwissenschaftliche und demokratisierte Umweltüberwachung.

\newpage

\section*{Anwendungsszenarien und Flexibilität}

Die Architektur der HydroNodeStation01 ist bewusst modular und erweiterbar gestaltet, um diverse Anwendungsszenarien zu unterstützen. Beispiele umfassen:

\begin{itemize}
  \item \textbf{Landwirtschaftliche Feldmessungen:} Erfassung von Bodenfeuchte, Bodenwärmestrom und Umgebungsparametern zur Optimierung von Bewässerung und Schädlingsbekämpfung in Regionen mit schlechter Netzabdeckung.
  
  \item \textbf{Wissenschaftliche Langzeitstudien:} Autonome Umweltüberwachung in Waldökosystemen, an Gewässern oder in Moorgebieten, wo eine manuelle Datenerfassung kaum praktikabel ist.
  
  \item \textbf{Luftqualitätsmessnetzwerke:} Dichte, partizipative Erfassung von Feinstaub und Gaskonzentrationen in städtischen und ländlichen Räumen ohne externe Infrastrukturabhängigkeit.
  
  \item \textbf{Früherkennung von Naturkatastrophen:} Dezentralisierte Sensornetzwerke zur Überwachung von Hochwasser, Erosion oder Rutschungsgefahren.
  
  \item \textbf{Bürgerwissenschaftliche Projekte:} Niedrigschwellige Partizipation von Bürgern und Schulen bei der Erfassung von Umweltdaten ohne hohe technische oder finanzielle Hürden.
\end{itemize}

Die offene I2C-Architektur ermöglicht es, beliebige kompatible Sensoren hinzuzufügen oder zu ersetzen. Dies macht das System adaptiv für zukünftige Messaufgaben und skalierbar für unterschiedliche Budgets und Anforderungen.

\section*{Struktur dieser Arbeit}

Die vorliegende Dokumentation beschreibt systematisch die Planung, Umsetzung und Validierung des Gesamtsystems. Neben den technischen Entscheidungen werden Aufwand, aufgetretene Probleme und gewählte Lösungsansätze transparent dargestellt. Dies ermöglicht eine vollständige Nachvollziehbarkeit des Projektablaufs und bildet die Grundlage für zukünftige Erweiterungen und Reproduktionen.
